\begin{remark}[The caveat that was missing, and what it is worth]\label{rem:type-caveat}
\texttt{METRIC\_TYPE\_UNKNOWN} is new to this table and is the reason to report caveats at
all rather than to list them once in prose. An \texttt{avg}, \texttt{min}, \texttt{max} or
\texttt{count} rollup reduces the points Datadog stores, which are samples for a gauge and
per-interval deltas for a counter; the emitted \texttt{*\_over\_time} is exact for the
first and unrelated to the second, and a dashboard does not record which it is
(Remark~\ref{rem:gauge}). At $100\%$ of translated occurrences it is the largest
qualification here, and it was missing from this table until the soundness theorem of
Section~\ref{sec:kernel-soundness} was written down and its hypotheses compared against
what the emitter actually recorded. It reached $100\%$ later still: the quarter of
translations that use \texttt{.as\_count()} or \texttt{.as\_rate()} were held to
establish their own metric type until the modifiers were measured on a gauge and found to
be silent no-ops (Remark~\ref{rem:gauge}). $\kappa$ did not move when it was added --- the
expressions are the same ones --- which is exactly why it is worth having: an unrecorded
precondition costs nothing in any number the paper reports, right up to the point where
somebody relies on the translation.

How often it is \emph{violated} is now measured rather than left as a worry, and it is
the one place in this paper where a precondition turns out to fail at a rate worth
reporting. The repository the dashboards come from declares each metric's type in a
file the corpus fetch had passed over, which settles the question for $7{,}557$ of the
$8{,}444$ translated occurrences --- $89.5\%$. On $491$ of those the front end's
inference disagrees with the declaration: $6.5\%$ of the decidable set and $5.6\%$ of
the corpus, split $242$ counters translated as gauges and $249$ gauges translated as
counters (Remark~\ref{rem:type-measured}). So $\kappa$ counts those $491$ among its
translations and they are, on the vendor's own declaration, wrong. $\kappa$ is a
coverage measure and was never a correctness measure; this is the first number in the
paper that says how far apart the two are.

Adding it also falsifies something an earlier draft of this section claimed. That draft
said $116$ occurrences carried \texttt{ROLLUP\_ALIGNMENT} alone and were the only
translations resting on nothing but the grid. With this code in place the true
figure is \emph{none}: every one of those $116$ is an \texttt{avg} rollup and now carries
\texttt{METRIC\_TYPE\_UNKNOWN} too. No translation in this corpus is exact on the strength
of the evaluation grid alone, which is a stronger claim than the one it replaces and was
only invisible because a precondition was. The same correction then ran to its conclusion:
with the modifiers' evidential value withdrawn, no translation in this corpus is exact on
the strength of anything the query itself contains.
\end{remark}
